% Part: incompleteness
% Chapter: representability-in-q
% Section: sigma1-completeness

\documentclass[../../../include/open-logic-section]{subfiles}

\begin{document}

\olfileid[tr]{inc}{inp}{s1c}
\olsection{\texorpdfstring{$\Sigma_1$}{Sigma-1} tamlığı}

$\Th{Q}$ ve onun tutarlı, aksiyomlaştırılabilir genişlemeleri eksik olsa da
$\Th{Q}$'nun sayısallar hakkında birçok temel olguyu ispatladığını gördük.
Aslında bu sonuç epey genişletilebilir. $\Th{Q}$ içinde ispatlanabilenlerin
kapsamını anlamak için $\Delta_0$, $\Sigma_1$ ve $\Pi_1$ !!{formula}s{}i
kavramlarını tanıtıyoruz. Kabaca bir $\Sigma_1$ !!{formula}{}ü, $!B$ yalnızca
önerme bağlaçları ve sınırlı niceleyiciler kullanılarak kurulmak üzere
$\lexists[x][!B(x)]$ biçimindedir. $!A$, $\Struct{N}$ içinde doğru olan bir
$\Sigma_1$ !!{sentence}{}si ise $\Th{Q} \Proves !A$ olduğunu göstereceğiz
(\olref{thm:sigma1-completeness}).

\begin{defn}
\ollabel{defn:bd-quant}
Bir \emph{sınırlı varlıksal !!{formula}}, $t$ herhangi bir terim olmak üzere
$\lexists[x][(x<t\land !A(x))]$ biçimindedir; bunu uzlaşımsal olarak
$\bexists{x<t}{!A(x)}$ diye yazarız.
%
Bir \emph{sınırlı evrensel !!{formula}}, $t$ herhangi bir terim olmak üzere
$\lforall[x][(x<t\lif !A(x))]$ biçimindedir; bunu uzlaşımsal olarak
$\bforall{x<t}{!A(x)}$ diye yazarız.
\end{defn}

\begin{defn}
\ollabel{defn:delta0-sigma1-pi1-frm}
Bir !!{formula} $!B$, atomik !!{formula}s{}den yalnızca önerme bağlaçları ve
sınırlı niceleme kullanılarak kurulmuşsa $\Delta_0$'dır.
%
Bir !!{formula} $!A$, $!B$ bir $\Delta_0$ formülü olmak üzere
$!A\ident\lexists[x][!B(x)]$ ise $\Sigma_1$'dir.
%
Bir !!{formula} $!A$, $!B$ bir $\Delta_0$ formülü olmak üzere
$!A\ident\lforall[x][!B(x)]$ ise $\Pi_1$'dir.
\end{defn}

\begin{lem}
\ollabel{lem:q-proves-clterm-id} $t$, $\Value{t}{N}=n$ olan kapalı bir terim
olsun. O zaman $\Th{Q} \Proves \eq[t][\num n]$ olur.
\end{lem}

\begin{proof}
Bunu $t$'nin karmaşıklığı üzerinde tümevarımla ispatlarız. Başlangıç
durumunda $\Value{\Obj 0}{N}=0$ olur ve $\num 0\ident\Obj 0$ olduğundan
$\Th{Q} \Proves \eq[\Obj 0][\num 0]$ elde edilir.
%
Tümevarım durumu için $t_1$ ve $t_2$,
$\Value{t_1}{N}=n_1$, $\Value{t_2}{N}=n_2$,
$\Th{Q} \Proves \eq[t_1][\num n_1]$ ve
$\Th{Q} \Proves \eq[t_2][\num n_2]$ olan terimler olsun.

O zaman $\Value{(t_1')}{N}=n_1+1$ olur. Tümevarım varsayımına ve
!!{formula} $\eq[\num{n_1}'][\num{n_1}']$'ye birinci derece özdeşlik
kurallarını uygulayarak $\Th{Q} \Proves \eq[t_1'][{\num n_1}']$ elde
ederiz; sayısalların tanımı uyarınca da
$\Th{Q} \Proves \eq[t_1'][\num{n_1+1}]$ olur.

Toplamlar için
\[
      \Value{(t_1 + t_2)}{N}
    = \Value{t_1}{N} + \Value{t_2}{N}
    = n_1 + n_2.
\]
Tümevarım varsayımı ve özdeşlik kuralları uyarınca
$\Th{Q} \Proves \eq[t_1+t_2][\num{n_1}+t_2]$ ve özdeşlik kurallarının
ikinci kez uygulanmasıyla
$\Th{Q} \Proves \eq[t_1+t_2][\num{n_1}+\num{n_2}]$ olur.
\olref[inc][req][bre]{lem:q-proves-add} uyarınca
$\Th{Q} \Proves \eq[\num{n_1}+\num{n_2}][\num{n_1+n_2}]$ olduğundan
$\Th{Q} \Proves \eq[t_1+t_2][\num{n_1+n_2}]$ elde edilir.

Benzer akıl yürütme \olref[inc][req][bre]{lem:q-proves-mult} kullanılarak
$\times$ için de işler.
%
Bu, aritmetiğin kapalı terimlerinin bütününü tükettiğinden
$\Value{t}{N}=n$ olan her kapalı $t$ terimi için
$\Th{Q} \Proves \eq[t][\num n]$ sonucuna ulaşırız.
\end{proof}

\begin{prob}
\olref{lem:q-proves-clterm-id}'in $\times$ içeren bölümünü ayrıntılı olarak
ispatlayın.
\end{prob}

\begin{lem}
\ollabel{lem:atomic-completeness}
$t_1$ ile $t_2$ kapalı terimler olsun. O zaman
\begin{enumerate}
\item $\Value{t_1}{N}=\Value{t_2}{N}$ ise
    $\Th{Q} \Proves \eq[t_1][t_2]$ olur.
\item $\Value{t_1}{N}\neq\Value{t_2}{N}$ ise
    $\Th{Q} \Proves \eq/[t_1][t_2]$ olur.
\item $\Value{t_1}{N}<\Value{t_2}{N}$ ise
    $\Th{Q} \Proves t_1<t_2$ olur.
\item $\Value{t_2}{N}\leq\Value{t_1}{N}$ ise
    $\Th{Q} \Proves \lnot(t_1<t_2)$ olur.
\end{enumerate}
\end{lem}

\begin{proof}
$t_1$ ve $t_2$ terimleri verildiğinde $n=\Value{t_1}{N}$ ve
$m=\Value{t_2}{N}$ değerlerini sabitleyelim.

$!A\ident t_1=t_2$ olduğunu varsayın. \olref{lem:q-proves-clterm-id}
uyarınca $\Th{Q} \Proves \eq[t_1][\num n]$ ve
$\Th{Q} \Proves \eq[t_2][\num n]$ olur. $n=m$ ise
$\Th{Q} \Proves \eq[\num n][\num m]$ ve dolayısıyla özdeşliğin geçişliliği
uyarınca $\Th{Q} \Proves \eq[t_1][t_2]$ olur. $n\neq m$ ise
$\Th{Q} \Proves \eq/[\num n][\num m]$ ve yine özdeşliğin geçişliliği
uyarınca $\Th{Q} \Proves \eq/[t_1][t_2]$ olur.

Şimdi $!A\ident t_1<t_2$ olsun. Her iki durumda da bütün $x,y$ için
$x<y\liff\lexists[z][\eq[z'+x][y]]$ diyen $!Q_8$ aksiyomuna dayanırız.

$\Sat{N}{t_1<t_2}$ olduğunu varsayın. O zaman $n+k+1=m$ olacak bir
$k\in\Nat$ vardır. \olref{lem:q-proves-clterm-id} uyarınca
$\Th{Q} \Proves \eq[t_1][\num n]$ ve
$\Th{Q} \Proves \eq[t_2][\num m]$ olur; bu lemmanın ilk bölümü uyarınca da
$\Th{Q} \Proves \eq[\num n+{\num k}'][\num m]$ olur. Özdeşliğin
geçişliliğinden $\Th{Q} \Proves \eq[{\num k}'+t_1][t_2]$ ve dolayısıyla
$\Th{Q} \Proves \lexists[z][\eq[z'+t_1][t_2]]$ çıkar. $!Q_8$'in sağdan
sola yönü uyarınca $\Th{Q} \Proves t_1<t_2$ olur.

Bunun yerine $\Sat/{N}{t_1<t_2}$, yani $m\leq n$ olduğunu varsayın.
%
$\Th{Q}$ içinde çalışıp $t_1<t_2$ olduğunu varsayıyoruz. $!Q_8$'in soldan
sağa yönü uyarınca $\eq[z'+t_1][t_2]$ olacak bir $z$ vardır.
$\Th{Q} \Proves \eq[t_1][\num n]$ ve
$\Th{Q} \Proves \eq[t_2][\num m]$ olduğundan
$\eq[z'+\num n][\num m]$ olur.
%
$!Q_5$ kullanılarak $m$ üzerinde dışsal bir tümevarımla
$\eq[z'+\num{n-m}][\Obj 0]$ elde edilir.
% [tr source emendation] Upstream prints a negated equality and cites Q_3;
% the derived equality contradicts the no-successor-is-zero axiom Q_2.
$m=n$ ise $\eq[z'][\Obj 0]$ olur ve bu $!Q_2$ aracılığıyla çelişki verir.
$m<n$ ise yine $!Q_5$ uyarınca
$\eq[(z'+\num{n-m-1})'][\Obj 0]$ olur ve bu da $!Q_2$ aracılığıyla çelişki
verir. Dolayısıyla $\Th{Q} \Proves \lnot(t_1<t_2)$ olur.
\end{proof}

\begin{lem}
\ollabel{lem:bounded-quant-equiv}
$!A$ !!a{formula}, $t$ kapalı bir terim ve $k=\Value{t}{N}$ olsun. O zaman
\begin{enumerate}
\item $\Th{Q} \Proves \bforall{x<t}{!A(x)}$ ancak ve ancak $\Th{Q} \Proves
    !A(\num 0) \land \dots \land !A(\num{k-1})$ olur.
\item $\Th{Q} \Proves \bexists{x<t}{!A(x)}$ ancak ve ancak $\Th{Q} \Proves
    !A(\num 0) \lor \dots \lor !A(\num{k-1})$ olur.
\end{enumerate}
\end{lem}

\begin{proof}
    Sınırlı evrensel niceleyici durumunu ispatlayalım.
    $\Value{t}{N}=0$ ise denkliğin sol tarafı $\Th{Q}$ içinde
    ispatlanabilir; çünkü \olref[inc][req][min]{lem:less-zero} uyarınca
    hiçbir $x<\num 0$ yoktur. Benzer biçimde boş bir birleşimi doğrudan
    $\ltrue$ sayabiliriz; bu da $\Th{Q}$ içinde ispatlanabilir.
    % [tr source emendation] Upstream says “empty disjunction” here, but the
    % bounded-universal expansion is the empty conjunction, with value true.
    %
    Bu nedenle bir $k$ doğal sayısı için $\Value{t}{N}=k+1$ olduğunu
    varsayalım. \olref{lem:q-proves-clterm-id} uyarınca
    $\bforall{x<\num{k+1}}{!A(x)}$ biçimindeki !!a{formula} ile
    çalıştığımızı varsayabiliriz.

    $\Th{Q} \Proves \bforall{x<\num{k+1}}{!A(x)}$ olduğunu varsayın ve
    $n\leq k$ olsun. \olref{lem:atomic-completeness} uyarınca
    $\Th{Q} \Proves \num n<\num{k+1}$ olduğundan mantık yoluyla
    $\Th{Q} \Proves !A(\num n)$ çıkar. Bu olguyu her $n\leq k$ için
    $k+1$ kez uygulayarak, istendiği gibi
    $\Th{Q} \Proves !A(\num 0)\land\dots\land !A(\num k)$ elde ederiz.

    Öteki yön için $\Th{Q} \Proves
    !A(\num 0)\land\dots\land !A(\num k)$ olduğunu varsayın. $\Th{Q}$
    içinde çalışarak $x<\num{k+1}$ olduğunu varsayın.
    \olref[inc][req][min]{lem:less-nsucc} uyarınca
    $x=\num 0\lor\dots\lor x=\num k$ olur; böylece mantık yoluyla $!A(x)$
    ve dolayısıyla evrensel iddia $\bforall{x<\num{k+1}}{!A(x)}$ çıkar.

    Sınırlı varlıksal olarak nicelenmiş !!{formula}s için denkliğin ispatı
    benzerdir.
\end{proof}

\begin{prob}
\olref{lem:bounded-quant-equiv}'deki varlıksal durumu ayrıntılı olarak
ispatlayın.
\end{prob}

\begin{lem}
\ollabel{lem:delta0-completeness}
$!A$, $\Struct{N}$ içinde doğru olan bir $\Delta_0$ !!{sentence}{}si ise
$\Th{Q} \Proves !A$ olur.
\end{lem}

\begin{proof}
Bunu !!{formula} karmaşıklığı üzerinde tümevarımla ispatlarız.
%
Başlangıç durumu \olref{lem:atomic-completeness} tarafından verilir;
dolayısıyla tümevarım adımına geçeriz. Yalınlık için değilleme durumunu,
değillemenin uygulandığı !!{formula}{}ün yapısına göre alt durumlara ayırırız.

\begin{enumerate}
\item $(!A\land !B)$, $\Struct{N}$ içinde doğru olsun; böylece $!A$ ile
$!B$, $\Struct{N}$ içinde doğrudur. Tümevarım varsayımı uyarınca
$\Th{Q} \Proves !A$ ve $\Th{Q} \Proves !B$ olur; dolayısıyla mantık
uyarınca $\Th{Q} \Proves (!A\land !B)$ elde edilir.
%
\item $\lnot(!A\land !B)$, $\Struct{N}$ içinde doğru olsun; böylece
$\lnot !A$ veya $\lnot !B$, $\Struct{N}$ içinde doğrudur. Genelliği
yitirmeden ilkinin doğru olduğunu varsayın. Tümevarım varsayımı uyarınca
$\Th{Q} \Proves \lnot !A$ ve dolayısıyla mantık uyarınca
$\Th{Q} \Proves \lnot(!A\land !B)$ olur.
%
\item $(!A\lor !B)$, $\Struct{N}$ içinde doğru olsun; böylece ya $!A$,
$\Struct{N}$ içinde ya da $!B$, $\Struct{N}$ içinde doğrudur. Genelliği yitirmeden ilkinin geçerli
olduğunu varsayın. Tümevarım varsayımı uyarınca $\Th{Q} \Proves !A$ ve
dolayısıyla mantık uyarınca $\Th{Q} \Proves (!A\lor !B)$ olur.
%
\item $\lnot(!A\lor !B)$, $\Struct{N}$ içinde doğru olsun; böylece
$\lnot !A$ ile $\lnot !B$, $\Struct{N}$ içinde doğrudur. Tümevarım
varsayımı uyarınca $\Th{Q} \Proves \lnot !A$ ve
$\Th{Q} \Proves \lnot !B$ olur. Bunun sonucu olarak mantık uyarınca
$\Th{Q} \Proves \lnot(!A\lor !B)$ elde edilir.
%
\item $\bforall{x<t}{!A(x)}$, $\Struct{N}$ içinde doğru olsun; burada $t$
kapalı bir terim ve $k=\Value{t}{N}$'dir. Her $n<\Value{t}{N}$ için
$!A(\num n)$, $\Struct{N}$ içinde doğruysa tümevarım varsayımı ve mantık
uyarınca $\Th{Q} \Proves !A(\num 0)\land\dots\land !A(\num{k-1})$ olur.
\olref{lem:bounded-quant-equiv} uyarınca
$\Th{Q} \Proves \bforall{x<t}{!A(x)}$ çıkar.
%
\item $\bexists{x<t}{!A(x)}$ biçiminde !!a{sentence} bulunan sınırlı
varlıksal niceleyici durumu, sınırlı evrensel niceleyici durumuna benzerdir.
%
\item $\lnot\bforall{x<t}{!A(x)}$, $\Struct{N}$ içinde doğru olsun; burada
$t$ kapalı bir terimdir. Bu !!{sentence}, $\Th{Q}$ içinde denkliği
türetilebilen $\bexists{x<t}{\lnot !A(x)}$ !!{sentence}{}sine denktir;
dolayısıyla sınırlı varlıksal niceleyiciler için kullanılan akıl yürütmeyi
uygulayabiliriz.
%
\item Benzer biçimde $\lnot\bexists{x<t}!A(x)$, $\Struct{N}$ içinde doğru
olsun; burada $t$ kapalı bir terimdir. Bu !!{sentence}, $\Th{Q}$ içinde
$\bforall{x<t}{\lnot !A(x)}$'e denktir; dolayısıyla sınırlı evrensel
niceleyiciler için kullanılan akıl yürütmeyi uygulayabiliriz.
%
\item Son olarak $\lnot !A$, $\Struct{N}$ içinde doğru olsun. Geriye yalnızca
$!A$'nın atomik olduğu ve bir $\Delta_0$ !!{sentence}{}si $!B$ için
$\lnot !A\ident\lnot\lnot !B$ olduğu durumlar kalır. $!A$ atomikse
\olref{lem:atomic-completeness} uyarınca $\Th{Q} \Proves \lnot !A$ olur.
$\lnot !A\ident\lnot\lnot !B$ ise mantık uyarınca bu, $\Th{Q}$ içinde
$!B$'ye ispatlanabilir biçimde denktir; bu tümce $\Struct{N}$ içinde
doğrudur, çünkü $\lnot !A$, $\Struct{N}$ içinde doğrudur. Dolayısıyla tümevarım varsayımı
uyarınca $\Th{Q} \Proves \lnot !A$ elde ederiz.
\end{enumerate}
\end{proof}

\begin{prob}
\olref{lem:delta0-completeness}'deki varlıksal durumu ayrıntılı olarak
ispatlayın.
\end{prob}

\begin{thm}
\ollabel{thm:sigma1-completeness}
$!A$, $\Struct{N}$ içinde doğru olan bir $\Sigma_1$ !!{sentence}{}si ise
$\Th{Q} \Proves !A$ olur.
\end{thm}

\begin{proof}
$\lexists{x}!A(x)$, $\Struct{N}$ içinde doğru olan bir $\Sigma_1$
!!{sentence}{}si ise $s(x)=n$ ve $\Sat{N}{!A(x)}[s]$ olacak bir $n$ doğal
sayısı ve bir $s$ değişken ataması vardır. Sağlama bağıntısı hakkındaki
standart olgulardan $\Sat{N}{!A(\num n)}$ çıkar. Fakat $!A(\num n)$ bir
$\Delta_0$ !!{formula} olduğundan \olref{lem:delta0-completeness} uyarınca
$\Th{Q} \Proves !A(\num n)$ ve dolayısıyla mantık uyarınca
$\Th{Q} \Proves \lexists[x][!A(x)]$ elde ederiz.
\end{proof}

\end{document}
